\chapter{مركبات الكربونيل: الإضافات المحبة للنواة، والأسيتالات والحماية}\label{ch:b1:acetals-protection}

أذِب الغلوكوز في الماء تجد أن شيئًا منه لا يكاد يبقى على هيئة
الألدهيد المفتوح السلسلة المرسوم في الكتب: فالجزيء ينثني، وتضاف مجموعته
الهيدروكسيلية إلى مجموعته الألدهيدية، فتنغلق حلقة سداسية. والإضافة
نفسها لكحول إلى مجموعة كربونيل، مكررة مرة أخرى \omterm{def:b1:elementary-steps:catalyst}{بحفّاز}
حمضي، تعطي \omterm{def:b1:acetals-protection:hemiacetal}{الأسيتالات}، وهي مركبات ثابتة تجاه القواعد وتجاه أشد
\omterm{def:b1:electronic-effects:nucleophile}{المحبات للنواة} تفاعليةً. وهذا الثبات أداة: فالألدهيد الذي كان
سيهدمه \omterm{def:b1:organomagnesium:organometallic}{كاشف غرينيار} يمكن إخفاؤه على هيئة \omterm{def:b1:acetals-protection:hemiacetal}{أسيتال}، وحمله
عبر التفاعل، واسترجاعه في النهاية. يدرس هذا الفصل
مجموعة الكربونيل \omterm{def:b1:electronic-effects:nucleophile}{محبًا للإلكترونات}، وإضافات الماء والكحولات،
واستراتيجية \omterm{def:b1:acetals-protection:protecting-group}{الحماية}.

\begin{recall}
إضافات غرينيار إلى مركبات الكربونيل (\cref{ch:b1:organomagnesium})؛
وحاصل التفاعل $Q$ وتطور الجملة ما دام $Q \neq K$
(\cref{ch:b1:extent-q-and-k})؛ \omterm{def:b1:electronic-effects:curly-arrow}{والأسهم المنحنية}، \omterm{def:b1:electronic-effects:nucleophile}{والمحبات للنواة}
\omterm{def:b1:electronic-effects:nucleophile}{والمحبات للإلكترونات} (\cref{ch:b1:electronic-effects}). وقدّم المجلد المدرسي (الصف
الثاني عشر) فكرة \omterm{def:b1:acetals-protection:protecting-group}{حماية} مجموعة أثناء التخليق.
\end{recall}

\section{مجموعة الكربونيل محبًا للإلكترونات}

الرابطة C=O مستقطبة نحو الأكسجين (\cref{ch:b1:periodicity})،
\omterm{def:b1:lewis-resonance-vsepr:resonance}{وبنية الرنين} \ce{C+-O-} تضع شحنة موجبة كاملة على الكربون:
فالكربون \omterm{def:b1:electronic-effects:nucleophile}{محب للإلكترونات}، تهاجمه \omterm{def:b1:electronic-effects:nucleophile}{المحبات للنواة} عموديًا على
مستوي المجموعة. \omterm{def:b1:electronic-effects:nucleophile}{والمحبات للنواة} القوية (كواشف غرينيار، والهيدريد،
والسيانيد) تُضاف مباشرة. أما \omterm{def:b1:electronic-effects:nucleophile}{المحبات للنواة} الضعيفة --- الماء، والكحولات --- فتحتاج إلى عون.

\begin{definition}[تنشيط الطابع المحب للإلكترونات]\label{def:b1:acetals-protection:activation}
\emph{تنشيط الطابع المحب للإلكترونات}\index{تنشيط الطابع المحب للإلكترونات} لمجموعة
الكربونيل هو برتنتها (أو تناسقها مع \omterm{def:b1:lewis-resonance-vsepr:lewis-acid}{حمض لويس}) على
الأكسجين. فالكاتيون \ce{C=OH+}، الذي تتقاسم شحنته الموجبة ذرةُ الكربون
عبر \omterm{def:b1:lewis-resonance-vsepr:resonance}{بنية الرنين} \ce{C+-OH}، تهاجمه حتى \omterm{def:b1:electronic-effects:nucleophile}{المحبات للنواة}
الضعيفة.
\end{definition}

\begin{omfigure}
\setchemfig{atom sep=2.1em}
\schemestart
\chemfig{C(-[:150]R)(-[:210]H)=O}
\arrow{->[\ce{H+}]}
\chemfig{C(-[:150]R)(-[:210]H)=O^{+}-H}
\arrow{<->}
\chemfig{C^{+}(-[:150]R)(-[:210]H)-O-H}
\schemestop

{\small تنشيط ألدهيد بالحمض: برتنة الأكسجين
تنتج كاتيونًا تحمل بنية رنينه الشحنة على
الكربون.}
\end{omfigure}

\section{الهيدرات وأنصاف الأسيتالات}

يُضاف الماء إضافةً عكوسة إلى الألدهيدات والكيتونات، \ce{R2C=O + H2O <=>
R2C(OH)2}؛ والهيدرات ثانوية عادةً في حالة الكيتونات ويمكن أن تكون رئيسية في حالة
الميثانال. ويُضاف الكحول بالطريقة نفسها.

\begin{definition}[نصف الأسيتال والأسيتال]\label{def:b1:acetals-protection:hemiacetal}
\emph{نصف الأسيتال}\index{نصف أسيتال} مركب يحمل فيه كربون واحد
مجموعة OH ومجموعة OR معًا، ويتكوّن بإضافة كحول
إلى مجموعة كربونيل. و\emph{الأسيتال}\index{أسيتال} مركب يحمل فيه
كربون واحد مجموعتي OR، ويتكوّن من نصف أسيتال وجزيء كحول
ثانٍ، مع فقد الماء.
\end{definition}

\omterm{def:b1:acetals-protection:hemiacetal}{أنصاف الأسيتالات} المفتوحة السلسلة ثانوية عادةً في توازنها مع
الألدهيد والكحول؛ أما حين تنتمي OH وC=O إلى الجزيء
نفسه ويمكن أن تنغلق حلقة خماسية أو سداسية، فيسود \omterm{def:b1:acetals-protection:hemiacetal}{نصف الأسيتال}
الحلقي. وهذه حالة الغلوكوز، وحالة المركب الأبسط
5-هيدروكسي بنتانال.

\begin{omfigure}
\setchemfig{atom sep=2em}
\schemestart
\chemfig{HO-[:30]-[:-30]-[:30]-[:-30]-[:30]C(=[:90]O)-[:-30]H}
\arrow{<=>}
\chemfig{*6(O-(-[:-30]OH)-----)}
\schemestop

\medskip

{\small 5-هيدروكسي بنتانال ونصف أسيتاله الحلقي، حلقة سداسية
تحتوي الأكسجين. لقد أُضيفت OH الكربون 5 إلى ألدهيد
الكربون 1؛ والشكل الحلقي هو السائد، كما في الغلوكوز.}
\end{omfigure}

\section{الأسيتالات}

\begin{proposition}[آلية تكوّن الأسيتال]\label{prop:b1:acetals-protection:acetal-mechanism}
في الوسط الحمضي، يعطي الألدهيد أو الكيتون مع كحول نصفَ \omterm{def:b1:acetals-protection:hemiacetal}{الأسيتال}، ثم
الأسيتالَ، عبر الخطوات: برتنة C=O؛ وإضافة الكحول؛
وفقد بروتون (\omterm{def:b1:acetals-protection:hemiacetal}{نصف الأسيتال})؛ وبرتنة OH؛ وفقد الماء،
فيتكوّن \omterm{def:b1:electronic-effects:intermediates}{كاتيون كربوني} تثبّته مجموعة OR الباقية (أيون
أوكسوكاربينيوم)؛ وإضافة كحول ثانٍ؛ وفقد بروتون. وكل خطوة
عكوسة.
\end{proposition}

\begin{proof}
كل خطوة انتقالُ بروتون حمضي--قاعدي أو إضافةُ (أو فقدُ)
\omterm{def:b1:electronic-effects:nucleophile}{محب للنواة} إلى (أو من) كربون \omterm{def:b1:electronic-effects:nucleophile}{محب للإلكترونات}، كما وُصف في
\cref{ch:b1:electronic-effects}. وفقد الماء ممكن لأن
الكاتيون المتكوّن، $\mathrm{R_2C^+{-}OR}$، له \omterm{def:b1:lewis-resonance-vsepr:resonance}{بنية رنين} $\mathrm{R_2C{=}O^+R}$
تكون فيها لكل ذرة ثمانية (\cref{prop:b1:electronic-effects:carbocation-order}).
ويُستهلك الحمض في عمليات البرتنة ويُستعاد في عمليات
نزع البروتون: فهو \omterm{def:b1:elementary-steps:catalyst}{حفّاز}.
\end{proof}

\begin{omfigure}
\setchemfig{atom sep=1.9em}
\schemestart
\chemfig{C(-[:150]R)(-[:210]H)=O}
\arrow{->[\ce{H+}][]}
\chemfig{C^{+}(-[:150]R)(-[:210]H)-OH}
\arrow{->[$\mathrm{R'OH}$][\ce{-H+}]}
\chemfig{C(-[:150]R)(-[:210]H)(-[2]OR')-OH}
\schemestop

\medskip
\schemestart
\chemfig{C(-[:150]R)(-[:210]H)(-[2]OR')-OH}
\arrow{->[\ce{H+}][\ce{-H2O}]}
\chemfig{C^{+}(-[:150]R)(-[:210]H)-OR'}
\arrow{->[$\mathrm{R'OH}$][\ce{-H+}]}
\chemfig{C(-[:150]R)(-[:210]H)(-[2]OR')-OR'}
\schemestop
\par\vspace{6pt}

{\small تكوّن \omterm{def:b1:acetals-protection:hemiacetal}{الأسيتال}، على سطرين: أولًا \omterm{def:b1:acetals-protection:hemiacetal}{نصف الأسيتال} (التنشيط،
وإضافة الكحول، وفقد بروتون)، ثم \omterm{def:b1:acetals-protection:hemiacetal}{الأسيتال} (برتنة
OH، وفقد الماء نحو الكاتيون المثبَّت، وإضافة كحول
ثانٍ، وفقد بروتون).}
\end{omfigure}

للتفاعل الإجمالي، $\mathrm{RCHO} + 2\,\mathrm{R'OH} \rightleftharpoons \mathrm{RCH(OR')_2} + \ce{H2O}$،
ثابت توازن متواضع، والماء أحد نواتجه.

\begin{proposition}[دفع تكوين الأسيتال]\label{prop:b1:acetals-protection:dean-stark}
إذا أُزيل الماء المتكوّن باستمرار، استمر تكوين \omterm{def:b1:acetals-protection:hemiacetal}{الأسيتال}
حتى يُستنفد مركب الكربونيل (أو الكحول).
\end{proposition}

\begin{proof}
يبقى حاصل التفاعل $Q = [\text{أسيتال}][\ce{H2O}]/([\text{ألدهيد}]
[\mathrm{R'OH}]^2)$ أدنى من $K$ ما دام الماء يُزال: فحسب
\cref{ch:b1:extent-q-and-k} يواصل التفاعل تقدمه، مهما كانت
قيمة $K$. ومع ثنائي أول (الإيثان-1,2-ثنائي أول) يكون \omterm{def:b1:acetals-protection:hemiacetal}{الأسيتال} حلقيًا، ويكون
التوازن أيضًا أكثر ملاءمة، إذ يحل جزيء واحد من ثنائي الأول محل جزيئين من
الكحول.
\end{proof}

\begin{omfigure}
\begin{tikzpicture}[font=\small]
  \pic at (0,0) {heatingmantle};
  \pic at (0,0) {rbflask={0.4}{omDef!14}};
  \pic (d) at (0,3.0) {deanstark={0.35}{omDef!35}{orange!15}};
  \pic at (0,6.4) {condenser};
  \draw[->] (2.4,4.0) node[right, align=left] {\foreignlanguage{arabic}{المجمِّع: الماء، الطبقة السفلى،}\\لا يعود؛ أما التولوين\\فيفيض عائدًا إلى الدورق} -- (0.8,3.8);
  \draw[->] (-2.2,1.0) node[left, align=right] {\foreignlanguage{arabic}{ألدهيد، ثنائي أول، تولوين،}\\حفّاز حمضي، غليان} -- (-0.7,0.9);
  \draw[->] (1.8,8.0) node[right] {مكثف} -- (0.4,7.8);
\end{tikzpicture}

{\small مصيدة دين--ستارك. تتكاثف أبخرة التولوين والماء
وتسقط في المجمِّع المدرَّج؛ والماء، الأكثف وغير الممتزج مع
التولوين، يهبط ويبقى، بينما يفيض التولوين عائدًا إلى الدورق.
وحجم الماء المجمَّع يقيس تقدم التفاعل.}
\end{omfigure}

\begin{proposition}[ثبات الأسيتالات]\label{prop:b1:acetals-protection:acetal-stability}
\omterm{def:b1:acetals-protection:hemiacetal}{الأسيتالات} ثابتة في الأوساط المتعادلة والقاعدية وتجاه \omterm{def:b1:electronic-effects:nucleophile}{المحبات للنواة}،
وكواشف غرينيار وكواشف الهيدريد؛ ويحلمئها الحمض المائي فيعيدها
إلى مركب الكربونيل والكحول.
\end{proposition}

\begin{proof}
لا يحمل كربون \omterm{def:b1:acetals-protection:hemiacetal}{الأسيتال} أي \omterm{def:b1:electronic-effects:nucleophile}{مجموعة مغادرة} يمكن لقاعدة أو \omterm{def:b1:electronic-effects:nucleophile}{لمحب للنواة}
أن يطردها: فالأيون \ce{RO-}، وهو \omterm{def:b1:predominant-reaction:strong-acid}{قاعدة قوية}، لا يغادر، ولا تبقى C=O
لتُهاجَم. وفي الحمض المائي تجري الآلية السابقة بالعكس: إذ يجعل
الفائض الكبير من الماء $Q < K$ للتفاعل العكسي.
\end{proof}

\section{الحماية وإزالة الحماية}

\begin{definition}[مجموعة الحماية]\label{def:b1:acetals-protection:protecting-group}
\emph{مجموعة الحماية}\index{مجموعة حماية} مجموعة تُدخل
في جزيء لتحجب وظيفة مؤقتًا، كي لا تتفاعل
في خطوة تستهدف جزءًا آخر من الجزيء. ويسمى إدخالها
\emph{الحماية}\index{حماية}، ونزعها
\emph{إزالة الحماية}\index{إزالة الحماية}.
\end{definition}

\begin{method}[احمِ، ثم فاعِل، ثم أزل الحماية]\label{met:b1:acetals-protection:protect}
\begin{enumerate}
  \item عيّن الوظيفة التي ستتفاعل، أو ستهدم
        الكاشف، في الخطوة المخطط لها.
  \item اختر \omterm{def:b1:acetals-protection:protecting-group}{مجموعة حماية} ثابتة في تلك الخطوة وقابلة للإزالة في
        ظروف يتحملها باقي الجزيء (في حالة ألدهيد أو
        كيتون يواجه قواعد أو \omterm{def:b1:electronic-effects:nucleophile}{محبات للنواة} أو هيدريدات: \omterm{def:b1:acetals-protection:hemiacetal}{أسيتال} حلقي).
  \item احمِ؛ ونفّذ الخطوة المخطط لها؛ ثم أزل \omterm{def:b1:acetals-protection:protecting-group}{الحماية}.
  \item احسب الكلفة: خطوتان إضافيتان، لكل منهما مردودها.
\end{enumerate}
\end{method}

\begin{example}[كيتون يعترض الطريق]\label{ex:b1:acetals-protection:ketoester}
لإضافة \omterm{def:b1:organomagnesium:organometallic}{كاشف غرينيار} إلى المجموعة الألدهيدية في جزيء يحمل أيضًا
كيتونًا، لا يوجد اختيار بسيط \omterm{def:b1:acetals-protection:hemiacetal}{للأسيتال} (فالألدهيدات تكوّن \omterm{def:b1:acetals-protection:hemiacetal}{الأسيتالات}
بسهولة أكبر من الكيتونات، فتُحمى المجموعة الخطأ)؛ ويجب تغيير
ترتيب الخطوات أو الكواشف. \omterm{def:b1:acetals-protection:protecting-group}{فالحماية}
استراتيجية، لا رد فعل آلي.
\end{example}

\section{تمارين}

\begin{exercise}[$\star$]\label{exo:b1:acetals-protection:1}
عيّن ذرات الكربون التي هي كربون \omterm{def:b1:acetals-protection:hemiacetal}{نصف أسيتال}، أو \omterm{def:b1:acetals-protection:hemiacetal}{أسيتال}، أو هيدرات، أو إيثر أو كحول في:
\ce{CH3CH(OH)OCH3}، \ce{CH3CH(OCH3)2}، \ce{CH3CH(OH)2}، \ce{CH3OCH3}،
\ce{(CH3)2C(OCH2CH3)2}.
\end{exercise}

\begin{exercise}[$\star$]\label{exo:b1:acetals-protection:2}
اكتب المعادلة الإجمالية لتكوّن \omterm{def:b1:acetals-protection:hemiacetal}{أسيتال} البروبانال مع
الميثانول، \omterm{def:b1:acetals-protection:hemiacetal}{والأسيتال} الحلقي للبروبانون مع الإيثان-1,2-ثنائي أول.
\end{exercise}

\begin{exercise}[$\star$]\label{exo:b1:acetals-protection:3}
أي هذه الكواشف يترك \omterm{def:b1:acetals-protection:hemiacetal}{الأسيتال} دون تغيير: محلول هيدروكسيد
الصوديوم، وبروميد ميثيل المغنيزيوم، وحمض الكبريتيك المخفف في الماء، وبوروهيدريد
الصوديوم؟
\end{exercise}

\begin{exercise}[$\star$]\label{exo:b1:acetals-protection:4}
ارسم \omterm{def:b1:acetals-protection:hemiacetal}{نصف الأسيتال} الحلقي الذي يكوّنه 4-هيدروكسي بوتانال. ما حجم
الحلقة؟
\end{exercise}

\begin{exercise}[$\star\star$]\label{exo:b1:acetals-protection:5}
اكتب الآلية الكاملة لتكوّن ثنائي ميثيل \omterm{def:b1:acetals-protection:hemiacetal}{أسيتال} الإيثانال
بحفز حمضي، \omterm{def:b1:electronic-effects:curly-arrow}{بالأسهم المنحنية}، وبيّن أن الحمض
يُستعاد.
\end{exercise}

\begin{exercise}[$\star\star$]\label{exo:b1:acetals-protection:6}
اكتب آلية الحلمأة الحمضية للمركب
2,2-ثنائي ميثوكسي بروبان. لماذا تُستعمل كمية كبيرة من الماء؟
\end{exercise}

\begin{exercise}[$\star\star$]\label{exo:b1:acetals-protection:7}
يُتابَع تكوين \omterm{def:b1:acetals-protection:hemiacetal}{أسيتال} \qty{0.250}{mol} من كيتون بمصيدة
دين--ستارك. ما حجم الماء الذي ينبغي جمعه عند التحول
التام (الكثافة نحو \qty{1.0}{g/mL})؟ وبعد ساعة،
تجمّع \qty{3.2}{mL}: فما نسبة التحول؟
\end{exercise}

\begin{exercise}[$\star\star$]\label{exo:b1:acetals-protection:8}
فسّر لماذا يتكوّن \omterm{def:b1:acetals-protection:hemiacetal}{الأسيتال} الحلقي مع الإيثان-1,2-ثنائي أول بصورة أتم
من \omterm{def:b1:acetals-protection:hemiacetal}{الأسيتال} مع جزيئين من الميثانول، لمركب الكربونيل
نفسه.
\end{exercise}

\begin{exercise}[$\star\star$]\label{exo:b1:acetals-protection:9}
لماذا يجب نزع \omterm{def:b1:elementary-steps:catalyst}{الحفّاز} الحمضي (أو معادلته) قبل معالجة المركب
المحمي \omterm{def:b1:organomagnesium:organometallic}{بكاشف غرينيار}؟
\end{exercise}

\begin{exercise}[$\star\star\star$]\label{exo:b1:acetals-protection:10}
اقترح تخليقًا للمركب 5-هيدروكسي-5-ميثيل هكسان-2-ون،
\ce{(CH3)2C(OH)CH2CH2COCH3}، من 4-كلوروبوتان-2-ون والبروبانون، مستعملًا
\omterm{def:b1:acetals-protection:protecting-group}{حماية}. أعطِ كل خطوة، وكواشفها ودور كل منها.
\end{exercise}

\begin{exercise}[$\star\star\star$]\label{exo:b1:acetals-protection:11}
بيّن بحاصل التفاعل أن فائضًا كبيرًا من الميثانول يدفع أيضًا
تكوين \omterm{def:b1:acetals-protection:hemiacetal}{أسيتال} الألدهيد، وقارن هذه الطريقة
بإزالة الماء.
\end{exercise}

\begin{exercise}[$\star\star\star$]\label{exo:b1:acetals-protection:12}
فسّر لماذا يُهاجَم \omterm{def:b1:acetals-protection:hemiacetal}{نصف الأسيتال} الحلقي للغلوكوز، في ميثانول
حمضي، فيعطي أسيتالًا ميثيليًا (غليكوزيدًا ميثيليًا) بينما لا تتحول مجموعات OH
الأخرى في الغلوكوز إلى إيثرات.
\end{exercise}

\section{مسألة: كاشف غرينيار يحمل ألدهيدًا}

\begin{problem}[{مسألة نهاية الأسبوع --- لماذا لا يستطيع 4-بروموبوتانال أن يعطي
كاشف غرينيار، وحمايته على هيئة أسيتال حلقي بمصيدة دين--ستارك،
وإضافة غرينيار إلى البروبانون، وإزالة الحماية،
وحجم الماء المجمَّع}]
\label{pb:b1:acetals-protection:1}
يحتاج كيميائي إلى 5-هيدروكسي-5-ميثيل هكسانال، \ce{(CH3)2C(OH)CH2CH2CH2CHO}،
ويخطط لتحضيره من 4-بروموبوتانال، \ce{BrCH2CH2CH2CHO}، والبروبانون.
ويبدأ من \qty{15.09}{g} من 4-بروموبوتانال. الكتل المولية
(\unit{g/mol}): \ce{C4H7BrO} 150.9، \ce{C2H6O2} 62.0، \ce{C6H11BrO2} 194.9،
\ce{C3H6O} 58.0، \ce{C9H18O3} 174.0، \ce{C7H14O2} 130.0، \ce{H2O} 18.0؛
كثافة الماء \qty{0.997}{g/mL}.
% ledger: rho:water

\medskip
\noindent\textbf{الجزء الأول --- المشكلة.}
\begin{enumerate}
  \item اكتب \omterm{def:b1:organomagnesium:organometallic}{كاشف غرينيار} الذي كان سيعطيه 4-بروموبوتانال.
  \item لماذا لا يمكن أن يوجد ذلك الكاشف؟ اكتب التفاعل الذي
        يهدمه.
  \item أي وظيفة يجب حمايتها، وضد أي شيء؟
  \item لماذا كان \omterm{def:b1:acetals-protection:hemiacetal}{الأسيتال} ملائمًا؟
  \item لماذا يُفضَّل \omterm{def:b1:acetals-protection:hemiacetal}{أسيتال} حلقي (مع الإيثان-1,2-ثنائي أول)؟
  \item احسب كمية 4-بروموبوتانال.
  \item احسب كتلة ثنائي الأول لفائض قدره \qty{20}{\percent}.
\end{enumerate}

\medskip
\noindent\textbf{الجزء الثاني --- \omterm{def:b1:acetals-protection:protecting-group}{الحماية}.}
\begin{enumerate}[resume]
  \item اكتب معادلة تكوين \omterm{def:b1:acetals-protection:hemiacetal}{الأسيتال} مع الإيثان-1,2-ثنائي أول.
  \item ما دور \omterm{def:b1:elementary-steps:catalyst}{الحفّاز} الحمضي (حمض 4-ميثيل بنزين
        السلفونيك)؟
  \item صِف عمل مصيدة دين--ستارك.
  \item فسّر بالمقدارين $Q$ و$K$ لماذا تدفع إزالة الماء التفاعل.
  \item لماذا يعود التولوين، لا الماء، إلى الدورق؟
  \item كيف يعرف الكيميائي أن التفاعل اكتمل؟
  \item احسب كتلة البروميد المحمي المتوقعة عند التحول التام.
\end{enumerate}

\medskip
\noindent\textbf{الجزء الثالث --- خطوة غرينيار.}
\begin{enumerate}[resume]
  \item اكتب تكوّن \omterm{def:b1:organomagnesium:organometallic}{كاشف غرينيار} من البروميد
        المحمي.
  \item اكتب إضافته إلى البروبانون \omterm{def:b1:alcohol-activation:alkoxide}{والألكوكسيد} المتكوّن.
  \item لماذا يجب أن تكون المعالجة في هذه الخطوة قاعدية خفيفة أو متعادلة (مثل
        كلوريد الأمونيوم المائي) لا حمضية؟
  \item احسب كتلة الكحول المحمي \omterm{def:b1:extent-q-and-k:fractional-extent}{لمردود} قدره \qty{70}{\percent}
        في هذه الخطوة.
\end{enumerate}

\medskip
\noindent\textbf{الجزء الرابع --- \omterm{def:b1:acetals-protection:protecting-group}{إزالة الحماية}.}
\begin{enumerate}[resume]
  \item اكتب \omterm{def:b1:acetals-protection:protecting-group}{إزالة الحماية} في الحمض المائي.
  \item لماذا يصمد الكحول الثالثي أمام الحمض المائي الخفيف؟
  \item يمكن للناتج أن يغلق حلقة سداسية بإضافة مجموعة OH فيه إلى
        الألدهيد. ارسم \omterm{def:b1:acetals-protection:hemiacetal}{نصف الأسيتال} الحلقي ذاك.
  \item احسب كتلة الألدهيد الهيدروكسيلي \omterm{def:b1:extent-q-and-k:fractional-extent}{بمردود} \qty{90}{\percent} في
        \omterm{def:b1:acetals-protection:protecting-group}{إزالة الحماية}.
  \item ما \omterm{def:b1:extent-q-and-k:fractional-extent}{المردود} الإجمالي الذي يبلغه الكيميائي انطلاقًا من 4-بروموبوتانال؟
  \item احسب كتلة الماء المتكوّن في خطوة \omterm{def:b1:acetals-protection:protecting-group}{الحماية} عند التحول
        التام.
  \item اذكر حجم الماء المجمَّع في مصيدة دين--ستارك عند التحول
        التام.
\end{enumerate}
\end{problem}
